%%%%%%%%%%%%%%%%%%%%%%%%%%%%%%%%%%%%%%%%%%%%%%%%%
%%%%%%%%%%%% chap: title %%%%%%%%%%%%%%%%%
%%%%%%%%%%%%%%%%%%%%%%%%%%%%%%%%%%%%%%%%%%%%%%%%%

\chapter{State of the art}\label{chapter:chap2}

The initial stage involves investigating existing work in the area of research to establish the foundation for the thesis.
\section{State of the Art in Phishing Honeypots and Behavioural Threat Intelligence}

\subsection{Phishing Attacks and Their Evolution}

The Phishing Industry has evolved from a few straightforward deceptive email messages to a sophisticated multi-channel and multi-layered architecture, including domain generation algorithms, CAPTCHA-based verification of a victim's identity, dynamically generating attack payloads from knowledge of the victims computer, and the ability to integrate Real-Time Command and Control \cite{Leita2008Phishing,Ramachandran2006FastFlux}. Today's phishing infrastructures are set up specifically to detect network intrusion detection (or IPSs) or decode security researchers' logs, thus the collection of Behavioural Intelligence (through Passive Monitoring) is becoming increasingly challenging \cite{Bringer2012HoneypotSurvey}.

\subsection{Traditional Phishing Detection Mechanisms}
Signature-based approaches identify malevolent URLs and domain names via the use of blacklists. Heuristic approaches evaluate webpage HTML components as well as URL character combinations/lettering. However, while these two approaches are effective in identifying previously known threats, they have a great deal of difficulty identifying unknown zero-day phishing threats and produce high amounts of false-negative results. While machine learning classifiers provide improved accurate detection of the threat, they are still limited to the static features of the URL that have been identified during the extraction process and do not provide any insights into what actions the victim will take after clicking through to the compromised website \cite{Bailey2009Botnets,Bringer2012HoneypotSurvey}.

\subsection{Honeypots in Cybersecurity}

Cybersecurity Experiences are considered to be low-interaction honeypots, high-interaction honeypots, and medium-interaction honeypots \cite{Seifert2006Taxonomy}. A low-interaction honeypot provides an environment that appears to be limited in its service capability while a high-interaction honeypot allows the individual a full exchange with the legitimate operating system and services. In most modern-day phishing research, the low-interaction honeypot is the predominant solution, as it most closely mimics the limited interaction that typically occurs during a phishing attack \cite{Artail2006Distributed}. As such, only superficial characteristics of the phishing campaign are captured; very little of the behavioural development of these campaigns is recorded \cite{Holz2006AttackPatterns}.

\subsection{Phishing-Specific Honeypots}

Current phishing honeypots are typically an exact duplicate of the legitimate login pages for popular services (i.e., email services and services) \cite{Provos2004ClientHoneypots}. The majority of these services will collect information about the user’s credentials and IP address; other sites will simply gather the phishing kit being utilized. Unfortunately, very few of these types of phishing honeypots offer the same realistic redirection paths or user verification methods that are incorporated into the more sophisticated phishing infrastructure \cite{Leita2008Phishing}.

\subsection{Behavioural Threat Intelligence}

Behavioral threat intelligence focuses on gaining understanding of what attackers actually do (and not just indicators) and therefore includes aspects of attacker activity such as navigation patterns, timing behaviours, reused infrastructure, and adaptive responses when detected \cite{Song2005HoneyStat,Holz2006AttackPatterns}. Therefore, it is key to creating predictive security systems that predict the attacker's actions before the attacker conducts the attacker's actions. Unfortunately, most intelligence feeds today are still based on passive detection rather than an active approach that extracts the behaviour of attackers \cite{Bringer2012HoneypotSurvey}.

\subsection{Automated Attacker Simulation}

Automated simulations are becoming more common as a method of stress-testing security controls. In the phishing space, automated tools like headless browsers or scripted bots allow for mass interactions with malicious actors \cite{Kreibich2005Replay,Alata2006Honeynet}. While the ability to conduct these types of studies has great potential, there have been few studies to date that have integrated automated attackers into honeypot environments to generate continuous behavioural data \cite{Bringer2012HoneypotSurvey}.

\subsection{Comparative Analysis of Existing Research}

Existing research demonstrates strong progress in botnet tracking, traffic replay, high-interaction honeypots, and large-scale attack measurement \cite{Bailey2009Botnets,Rajab2007Botnet,Freiling2010Evasion}. However, these works primarily focus on malware propagation and network-level attacks rather than phishing-specific multi-stage behavioural analysis.

\subsection{Identified Research Gaps}
For the literature reviewed, these gaps appear to exist in the following areas:

1. No multi-stage phishing emulation has been found in the existing honeypot literature \cite{Bringer2012HoneypotSurvey}.
2. Very limited studies have integrated automated attacker simulations \cite{Kreibich2005Replay}.
3. Very limited studies have collected end-to-end behavioural indicators on attackers \cite{Song2005HoneyStat}.
4. There has been a minimal emphasis on analysing the adaptive responses of the attackers \cite{Freiling2010Evasion}.
5. There are weak links between the data generated in honeypots and the ability to take action based on that data \cite{Bringer2012HoneypotSurvey}.

The proposed research will directly address each of these gaps via an integrated, automated, behaviourally rich phishing deception platform.

\section{Identified Opportunities for Publishing}

The research proposed here on a multi-stage phishing decoy honeypot using automatic attacker simulation and behaviour based threat intelligence generation addresses a current and significant issue that affects all aspects of cybersecurity. As this research encompasses many aspects of an interdisciplinary approach, namely web security, deception technology, automated attacker simulation, and the analysis of behavioural data, the results of this project will be useful for dissemination in numerous reputable international scientific venues. In this section we will identify and provide the reasoning behind various peer-reviewed conference and journal publication opportunities closely aligned with the methodology and contributions of this study.

\subsection{Security and Cyber Defence Conferences}

International Conferences are arguably the best category of publication venues for this research; Many of the top International Conferences host papers on research associated with Hacking and Network Security. The Annual Conference on Computer Security Applications (ACSAC) would serve as another excellent option to publish this research. ACSAC has been a very successful peer-reviewed International conference, which has a strong focus on applied security research, such as; Intrusion Detection, Malware Analysis, Cyber Deception and Attack Defence systems. The experimental design of building a multi-stage Phishing Honeypot, incorporating automated simulation of Attacker Behaviour, aligns with the practical and deployable elements of the security solution that ACSAC is seeking when compiling and publishing submissions for the conference. Additionally, the Behavioural Threat Intelligence generated from the platform would provide additional foundations upon which to support future use of defence systems based on this behaviour analysis.

Another International Conference that is highly suitable for the publication of this research is the European Symposium on Research in Computer Security (ESORICS). ESORICS is a European based International Conference that publishes work on both theoretical and applied Security Research. The common topics of publications from ESORICS include Network Security, Web Security, Adversarial Behaviour Modelling and Attack Detection. Methodologies developed through this work, such as; Automatic Behavioural Data Acquisition and Multi-Stage Deception Architectures, fit well within the overall scope of ESORICS. Publication in ESORICS would expose and disseminate this research to the broader academic community, as well as contribute to the creation of new knowledge and understanding around Deception-Based Security Mechanisms.

The USENIX Security Symposium is an additional major opportunity for getting high Impact. This symposium is one of the top conferences in the area of computer security. There are many submissions to become part of this symposium and it has a very competitive process for selecting which papers will be accepted. It is very common to see papers covering subjects such as “honeypots”, “web based attacks, phishing” and “large scale behaviour of adversarial networks.” The USENIX Security Symposium is also a prime venue to present the findings from the research when it comes to using empirical results and new automated techniques of presenting them.

\subsection{Web, Network, and Application Security Conferences}

Phishing attacks depend on Internet Technology and the infrastructure of the Internet. There are several conferences related to both web and network security that offer a wealth of opportunities for publishing research on phishing attacks. A prime example is the ACM Conference on Data \& Application Security and Privacy (CODASPY), which is an excellent venue due to the conference's strong emphasis on application-layer security, data protection, and privacy-aware system design. The proposed phishing honeypot platform will interact directly with web-based credential harvesting mechanisms and will capture behaviour-related information associated with the breach of privacy in a controlled way, which would make it an excellent match for the conference to present the System's architecture and ethical data collection framework. RAID is yet another appealing option for a venue to publish research on phishing. As its name implies, RAID focuses on the detection, analysis, and mitigation of cyber attacks, including phishing, botnets, malware, and various other evasive \& advanced cyber attack strategies. The adaptive analysis of behaviour and automated simulation of an attacker align closely with the goals of RAID, with its paramount focus on modelling advanced attacks and validating the mechanisms of defense against them. The components of the behavioural analysis and the profiling of attacker interactions are ideal for a submission to RAID.


\subsection{Scientific Journals}

In addition to Conference Publications there are a number of reputable/journals that are good opportunities to publish an Extended Version of your Work. One journal in particular that is great for Cyber Security Research is the IEEE Transactions on Information Forensics and Security (IEEE TIFS). They frequently publish Articles about: Intrusion Detection; Analysis of Attack Behaviour; Cyber Deception; and Security Systems Based on Data. An Extended Version of the Behavioural Intelligence Framework along with Comprehensive Experimental Validation would likely be a Good Fit for Submission to This journal.
Another General Purpose Journal that you might consider is Computers and Security published by Elsevier. This Journal focuses on Applied Security Research, Detection of Attacks, and Security Systems Design. The Proposed Multi-Stage Phishing Honeypot System, Along with Its Automated Data Collection and Data Analysis Pipeline are Well Aligned with the Journals Focus on Practical Solutions for Security and Empirical Validation.
A Third Good Option for Publication is The Journal of Computer Virology and Hacking Techniques which publishes a great deal of Work Related to Malicious Infrastructure, Attack Models and using Defensive Experiments. The Automated Attacker Simulation and Phishing Behavior Profiling components of this Research are Well Within the Journal's Focus Area.

\subsection{Publication Strategy and Expected Impact}

This research publication strategy has been organized into phases. The first phase will contain a conference submission that emphasizes the system architecture, deploying methods, and early behaviour analysis results. This will expedite the dissemination of the primary technical findings and allow researchers in this field to provide input regarding the study’s methodology and initial findings. Afterward, an extended analysis article will be published through the journal, including the additional refined experimentation, longitudinal behaviour analysis, and advanced statistical evaluations into how attackers react to the system.

The publication of the results of this research will serve as a high-impact publication that maximizes the potential for academic and practical relevance of this work. The academic value of this research will provide new methods of evaluating phishing behaviours utilizing experimental techniques and create a framework for automated detection of phishing schemes. From an industrial and society perspective, this research creates findings that will directly benefit operational cyber security centres, platforms for quantitative threat intelligence, and anti-phishing systems. Therefore, this work will help shape how future cyber defences are developed by increasing the standard of cyber deception and behaviour-based threat intelligence.

